\subsection{希尔伯特空间的希尔伯特张量积}

定义 1. --- 设 \(E_{1}, \ldots, E_{n}\) 为希尔伯特空间。分离准希尔伯特空间 \(E_{1} \otimes_{2} \ldots \otimes_{2} E_{n}\) 的完备化称为诸 \(E_{i}\) 的希尔伯特张量积，记作 \(E_{1} \hat{\otimes}_{2} \ldots \hat{\otimes}_{2} E_{n}\)（或 \(\hat{\otimes}_{2} E_{i}\)）。

设 \(F_{1}, \ldots, F_{n}\) 为希尔伯特空间，且 \(u_{i} \in \mathcal{L}(E_{i}, F_{i})\)（\(1 \leqslant i \leqslant n\)）。于是连续线性映射 \(u_{1} \otimes \ldots \otimes u_{n}\) 可延拓为连续线性映射 \(u_{1} \hat{\otimes}_{2} \ldots \hat{\otimes}_{2} u_{n}\)（从 \(E_{1} \hat{\otimes}_{2} \ldots \hat{\otimes}_{2} E_{n}\) 到 \(F_{1} \hat{\otimes}_{2} \ldots \hat{\otimes}_{2} F_{n}\) 中）。我们有

\[
\| u _ {1} \hat {\otimes} _ {2} \dots \hat {\otimes} _ {2} u _ {n} \| = \| u _ {1} \| \dots \| u _ {n} \|\tag{10}
\]

这由 V, p.~27 的公式 (8) 得到。此外，若以 \(1_{E}\) 表示任意希尔伯特空间 E 的恒等映射，则有

\[
1 _ {\mathrm{E} _ {1}} \hat {\otimes} _ {2} \dots \hat {\otimes} _ {2} 1 _ {\mathrm{E} _ {n}} = 1 _ {\mathrm{E}} \quad \text { with } \quad \mathrm{E} = \mathrm{E} _ {1} \hat {\otimes} _ {2} \dots \hat {\otimes} _ {2} \mathrm{E} _ {n}.\tag{11}
\]

最后，若 \(\mathbf{G}_1, \ldots, \mathbf{G}_n\) 为希尔伯特空间且 \(v_i \in \mathcal{L}(\mathbf{F}_i; \mathbf{G}_i)\)（\(1 \leqslant i \leqslant n\)），则得到

\[
\left(v _ {1} \circ u _ {1}\right) \hat {\otimes} _ {2} \dots \hat {\otimes} _ {2} \left(v _ {n} \circ u _ {n}\right) = \left(v _ {1} \hat {\otimes} _ {2} \dots \hat {\otimes} _ {2} v _ {n}\right) \circ \left(u _ {1} \hat {\otimes} _ {2} \dots \hat {\otimes} _ {2} u _ {n}\right).\tag{12}
\]

与上面对准希尔伯特空间所述的内容相类比，希尔伯特张量积的«交换性»与«结合性»如何表述，留给读者自己去完成。

注. --- 设 \(E_{1}, \ldots, E_{n}\) 为分离的准希尔伯特空间，\(\hat{E}_{1}, \ldots, \hat{E}_{n}\) 为它们各自的完备化。于是 \(E_{1} \otimes_{2} \ldots \otimes_{2} E_{n}\) 是 \(\hat{E}_{1} \otimes_{2} \ldots \otimes_{2} \hat{E}_{n}\) 的准希尔伯特子空间。由于映射 \((x_{1}, \ldots, x_{n}) \mapsto x_{1} \otimes \ldots \otimes x_{n}\)（从 \(\hat{E}_{1} \times \cdots \times \hat{E}_{n}\) 到 \(\hat{E}_{1} \otimes_{2} \ldots \otimes_{2} \hat{E}_{n}\) 中）是连续的，故 \(E_{1} \otimes_{2} \ldots \otimes_{2} E_{n}\) 在 \(\hat{E}_{1} \otimes_{2} \ldots \otimes_{2} \hat{E}_{n}\) 中稠密。更有甚者，\(E_{1} \otimes_{2} \ldots \otimes_{2} E_{n}\) 的完备化恰好就是希尔伯特空间 \(\hat{E}_{1} \otimes_{2} \ldots \otimes_{2} \hat{E}_{n}\)。这一完备化有时就简记作 \(E_{1} \otimes_{2} \ldots \otimes_{2} E_{n}\)（或 \(\bigotimes_{1 \leq i \leq n} E_{i}\)）。

命题 3. --- 设 \(E_{1}, \ldots, E_{n}\) 为希尔伯特空间。假定对 \(1 \leqslant i \leqslant n\)，空间 \(E_{i}\) 是闭向量子空间族 \((\mathrm{E}_{i,\alpha})_{\alpha \in \mathbf{A}(i)}\) 的希尔伯特和。于是 \(E_{1} \otimes_{2} \ldots \otimes_{2} E_{n}\) 是子空间族 \(E_{1,\alpha_{1}} \hat{\otimes}_{2} \ldots \hat{\otimes}_{2} E_{n,\alpha_{n}}\) 的希尔伯特和，其中 \((\alpha_{1}, \ldots, \alpha_{n})\) 取遍 \(A(1) \times \cdots \times A(n)\)。

由 V, p.~27 的公式 (6)，子空间 \(\mathrm{E}_{1,\alpha_1}\hat{\otimes}_2\ldots \hat{\otimes}_2\mathrm{E}_{n,\alpha_n}\)（\(\mathrm{E}_1\hat{\otimes}_2\ldots \hat{\otimes}_2\mathrm{E}_n\) 中的）两两正交。对每个整数 \(i\)（介于 1 与 \(n\) 之间），集合 \(\bigcup_{\alpha \in \mathbf{A}(i)}\mathrm{E}_{i,\alpha}\) 在 \(\mathrm{E}_i\) 中是完全的，且多线性映射 \((x_{1},\dots,x_{n})\mapsto x_{1}\otimes \dots \otimes x_{n}\) 是连续的。由此推出诸子空间 \(\mathrm{E}_{1,\alpha_1}\hat{\otimes}_2\ldots \hat{\otimes}_2\mathrm{E}_{n,\alpha_n}\) 的并是完全的，从而命题 3 成立。

推论 1. --- 对 \(1 \leqslant i \leqslant n\)，设 \((e_{i,\alpha})_{\alpha \in \mathbf{A}(i)}\) 为 \(\mathbf{E}_i\) 的标准正交基。于是向量族 \(e_{1,\alpha_1} \otimes \ldots \otimes e_{n,\alpha_n}\)（当 \((\alpha_1, \ldots, \alpha_n)\) 取遍 \(\mathbf{A}(1) \times \cdots \times \mathbf{A}(n)\) 时）是 \(\mathbf{E}_1 \hat{\otimes}_2 \ldots \hat{\otimes}_2 \mathbf{E}_n\) 的标准正交基。

推论 2. --- 设 \(E_{1}\) 与 \(E_{2}\) 为两个希尔伯特空间，\((e_{i})_{i\in I}\) 为 \(E_{1}\) 的标准正交基。设 \((y_{i})_{i\in I}\) 为 \(E_{2}\) 中元素的一个族，满足 \(\sum_{i\in I}\|y_{i}\|^{2}<+\infty\)。于是族 \((e_{i}\otimes y_{i})_{i\in I}\) 在 \(E_{1}\hat{\otimes}_{2}E_{2}\) 中可和；此外，\(E_{1}\hat{\otimes}_{2}E_{2}\) 的每个元素都可以唯一地写成 \(\sum_{i\in I}e_{i}\otimes y_{i}\) 的形式，其中 \(\sum_{i\in I}\|y_{i}\|^{2}<+\infty\)。

设 \(F_{i}\) 为 \(E_{1}\) 中由 \(e_{i} (i \in I)\) 生成的直线。于是 \(E_{1}\) 是子空间族 \((F_{i})_{i \in I}\) 的希尔伯特和。由命题 3，空间 \(E_{1} \hat{\otimes} E_{2}\) 是子空间族 \((F_{i} \hat{\otimes}_{2} E_{2})_{i \in I}\) 的希尔伯特和，由此即得推论 2。

例. --- 1) 由推论 1，空间 \(\ell^2 (\mathrm{I})\hat{\otimes}_2\ell^2 (\mathrm{J})\) 典范同构于 \(\ell^2 (\mathrm{I}\times \mathrm{J})\)，张量积 \(\mathbf{x}\otimes \mathbf{y}\)（其中 \(\mathbf{x} = (x_i)_{i\in \mathrm{I}}\)，\(\mathbf{y} = (y_j)_{j\in \mathrm{J}}\)）可与族 \((x_iy_j)_{i\in \mathrm{I},j\in \mathrm{J}}\) 等同。类似地，由推论 2，\(\ell^2 (\mathrm{I})\hat{\otimes}_2\mathrm{E}\) 可与 \(\ell_{\mathrm{E}}^{2}(\mathrm{I})\) 等同，使得 \((x_{i})_{i\in \mathrm{I}}\otimes y = (x_{i}y)_{i\in \mathrm{I}}\)（对希尔伯特空间 E 中的每个 \(\mathbf{y}\)）。

* 2) 设 X 为分离的拓扑空间，\(\mu\) 为 X 上的一个正测度。设 E 为希尔伯特空间。我们可以用典范的方式把 \(L^2(X, \mu) \otimes_2 E\) 与 \(L_E^2(X, \mu)\) 等同起来：若 \(f\) 是 X 上平方可积数值函数 \(f\) 的类，且 \(a\) 属于 E，则 \(f \otimes a\) 就是取值于 E 中的函数 \(x \mapsto f(x).a\) 的类。

设 Y 为分离的拓扑空间，v 为 Y 上的一个正测度。以类似的方式，我们可以把希尔伯特空间 \(\mathrm{L}^{2}(\mathrm{X},\mu)\hat{\otimes}_{2}\mathrm{L}^{2}(\mathrm{Y},v)\) 与 \(\mathrm{L}^{2}(\mathrm{X}\times\mathrm{Y},\mu\otimes v)\) 等同起来；这时 \(\dot{f}\otimes\dot{g}\) 可与函数 \((x,y)\mapsto f(x)g(y)\)（在 \(X\times Y\) 上）的类等同。

